%=====================================================================
\section{Revision 7: corrections, supersessions and changelog (19 October 2026)}\label{app:rev7}
%=====================================================================
Revision~7 adds \S\ref{sec:rev7}. No earlier repository file is edited, renamed or deleted; flags on earlier pushed files are recorded in new repository note files only (\path{NOTE_Rev7_Flags_on_Earlier_Files_2026-10-19.md} and \path{NOTE_LABEL_CFP.md}). CORE\_FREEZE holds: the locked $H_s$ ($c_2=(-4,1)$) is unchanged and $H_s^{\rm axes}$ is a parallel object. No claim is made concerning HQCC termination, the twin prime conjecture or the Riemann hypothesis.
\begin{enumerate}[label={R7.\arabic*},leftmargin=3em]
\item \textbf{Correction (role of $\Delta$).} The earlier statement that $\Delta$ constrains the S$^2$-side expansions intrinsically is refined: only $(s-2)$ is intrinsic given $c_1$; $7s-32$ and the $\sqrt\Delta$ terms are tied to the locked $c_2=(-4,1)$ (Remark~\ref{rem:corr7}, Category~A).
\item \textbf{Superseded (Category~B $\to$ A).} The $k(k-1)$ count of extra $h$-relations (\S\ref{sec:rev6h}, status box; open problem~1 of \S\ref{sec:open6}) is proved for all $k\ge2$ (Proposition~K1). The resulting dimension bounds are proved upper bounds; their sharpness is open.
\item \textbf{Superseded.} The status of Proposition~\ref{prop:crt6}(v) (``only a partial conceptual reason'' for the periods $16$ and $96$) is replaced by Propositions~\ref{prop:evenN7} and~\ref{prop:n8dist7}: $N=2,4$ fully conceptual; $N=8$ invariants derived by hand and distinctness proved by a residue-level finite computation, with no scan over $s$.
\item \textbf{Superseded.} ``The $t^{e/2}$ coefficients at $x=\tfrac14,\tfrac16$ were not computed'' (status of Proposition~\ref{prop:crt6}; open problem~3 of \S\ref{sec:open6}): they are now given at $t^{11/2},t^{13/2}$ (Proposition~\ref{prop:x1416}), together with two vanishing results.
\item \textbf{Label.} The CRT factorisation statement (Proposition~\ref{prop:crt6}(ii)) is called the CRT Factorization Proposition (CFP). The letter label used for it in [S36] and in two script comments is retired; see \path{NOTE_LABEL_CFP.md}. The $k(k-1)$ result is Proposition~K1.
\item \textbf{Source corrections recorded.} [S36], G.2: the $(s-8)$-exponent bullet is garbled; the correct exponent is $\max(0,e+3-4j)$ ([S36], G.0). [S37], \S7: the stored \texttt{c7/verify.json} is empty, so its engine-agreement sentence is not used; the earlier ``direct'' block there ($t\ge0.01$) is not a valid test and is superseded by \texttt{verify\_direct.json} (Category~B).
\item \textbf{(C3) extended.} Category~A now also for $\Z/47,\Z/59,\Z/71$ and $\Z/17^2,\Z/19^2,\Z/23^2$ (Proposition~\ref{prop:weil7}); general $p$ still cited. $s=121$, $3481$ remain factor-dimension predictions.
\item \textbf{Added:} \S\ref{sec:rev7K1}--\S\ref{sec:open7}; Proposition~K1; Propositions~\ref{prop:delta7}--\ref{prop:weil7}; Theorem~\ref{thm:axes7}; Definition~\ref{def:axes7}; Remarks~\ref{rem:intrinsic7}, \ref{rem:corr7}; sources [S36], [S37].
\end{enumerate}
